% 核心结论：滞后负载与超前负载下串联阻抗压降的相量方向
% 图原型：纯几何／对比。物理网络按本题节点、端口和符号构建。
\documentclass[tikz,border=7pt]{standalone}
\usepackage{fontspec,xeCJK,amsmath,amssymb}
\IfFontExistsTF{Times New Roman}{\setmainfont{Times New Roman}}{\setmainfont{TeX Gyre Termes}}
\IfFontExistsTF{Heiti SC}{\setCJKmainfont[AutoFakeBold=1.5]{Heiti SC}}{\setCJKmainfont[AutoFakeBold=1.5]{FandolHei-Regular}}
\usetikzlibrary{arrows.meta,calc,angles,quotes}
\definecolor{ink}{HTML}{243745}
\definecolor{blue}{HTML}{275DAB}
\definecolor{green}{HTML}{177C63}
\definecolor{red}{HTML}{BD4935}
\definecolor{purple}{HTML}{715393}
\begin{document}
\begin{tikzpicture}[font=\fontsize{9.5}{12}\selectfont,text=ink,
 v/.style={line width=1.1pt,-{Stealth[length=2.4mm,width=1.5mm]}},
 axis/.style={ink!65,line width=.45pt,-{Stealth[length=1.7mm,width=1.1mm]}},
 aux/.style={ink!40,densely dashed,line width=.5pt},
 lab/.style={fill=white,inner sep=2pt},
 card/.style={draw=blue!25,fill=blue!3,rounded corners=3pt,text width=6cm,align=left,inner sep=10pt}]
\path[use as bounding box] (-8,-3.35) rectangle (8,3.65);
\node[font=\fontsize{10.5}{13}\selectfont\bfseries,text width=15cm,align=center] at (0.00000,3.22000) {M6-H36　滞后负载与超前负载下串联阻抗压降的相量方向};
\begin{scope}[shift={(-4,0)}]
\draw[axis] (-2.20000,-0.50000) -- (2.00000,-0.50000);
\draw[v,blue] (-2.20000,-0.50000) -- (0.20000,-0.50000);
\draw[v,green] (0.20000,-0.50000) -- (0.77600,-0.93200);
\draw[v,purple] (0.77600,-0.93200) -- (1.35200,-0.16400);
\draw[v,red] (-2.20000,-0.50000) -- (1.35200,-0.16400);
\node[text=blue] at (-1.00000,-0.85000) {$U_2^{\prime}$};
\node[text=green] at (0.65000,-1.30000) {$IR$};
\node[text=purple] at (1.80000,0.55000) {$jIX$};
\node[text=red] at (-0.70000,1.10000) {$\widetilde U_1$};
\node[] at (0.00000,2.20000) {滞后电流};
\node[] at (0.00000,-2.40000) {$\Delta U_{\parallel}\simeq I(R\cos\phi+X\sin\phi)$};
\end{scope}
\begin{scope}[shift={(4,0)}]
\draw[axis] (-2.20000,-0.50000) -- (2.00000,-0.50000);
\draw[v,blue] (-2.20000,-0.50000) -- (0.20000,-0.50000);
\draw[v,green] (0.20000,-0.50000) -- (0.77600,-0.06800);
\draw[v,purple] (0.77600,-0.06800) -- (0.20000,0.70000);
\draw[v,red] (-2.20000,-0.50000) -- (0.20000,0.70000);
\node[text=blue] at (-1.00000,-0.85000) {$U_2^{\prime}$};
\node[text=green] at (0.65000,-1.30000) {$IR$};
\node[text=purple] at (1.80000,0.55000) {$jIX$};
\node[text=red] at (-0.70000,1.10000) {$\widetilde U_1$};
\node[] at (0.00000,2.20000) {超前电流};
\node[] at (0.00000,-2.40000) {$\Delta U_{\parallel}\simeq I(R\cos\phi-X\sin\phi)$};
\end{scope}
\node[text width=15cm,align=center] at (0.00000,-3.00000) {比较同一电流大小下压降的实轴投影；近似纵向压降为零不表示精确电压模差严格为零。};

\end{tikzpicture}\end{document}